%%% ============================================================
%%% The Sinc^2 Zero Law for Squarefree Moduli
%%% B. R. Sanders (7Site LLC, Hot Springs AR), M. Gish, 2026
%%% DOI: 10.5281/zenodo.18852047
%%% Target venue: Integers -- Electronic Journal of
%%%               Combinatorial Number Theory
%%% Source: WP_SINC2_ZERO_LAW.md (re-scoped 2026-05-07 from prime
%%%         to squarefree modulus, per the Sprint 35 audit)
%%% Verification: proof_d25_loop_closure.py (primes 3..199, exact)
%%% MSC: 11A41, 11N05, 42A16
%%%
%%% NOTE: This manuscript replaces the 2026-04-18 prime-only draft
%%% which was pulled back on 2026-04-19 because the basic
%%% biconditional sinc^2(k/n) = 0 <=> n | k is uniform in n. The
%%% present draft re-scopes the result to squarefree b and routes
%%% the prime-specific content through the First-G companion paper
%%% (J04, Sanders + Gish 2026; companion submission to Integers).
%%% ============================================================

\documentclass[11pt,reqno]{amsart}

\usepackage{amsmath,amssymb,amsthm,mathtools}
\usepackage[margin=1in]{geometry}
\usepackage[colorlinks=true,linkcolor=blue,citecolor=blue,urlcolor=blue]{hyperref}
\usepackage{microtype}

%% -------- theorem environments --------
\theoremstyle{plain}
\newtheorem{theorem}{Theorem}
\newtheorem{corollary}[theorem]{Corollary}
\newtheorem{lemma}[theorem]{Lemma}
\newtheorem{proposition}[theorem]{Proposition}

\theoremstyle{remark}
\newtheorem*{remark}{Remark}

%% -------- macros --------
\DeclareMathOperator{\sinc}{sinc}
\DeclareMathOperator{\spf}{spf}
\newcommand{\Z}{\mathbb{Z}}
\newcommand{\N}{\mathbb{N}}
\newcommand{\Q}{\mathbb{Q}}
\newcommand{\Prm}{\mathbb{P}}

%% -------- title matter --------
\title[The sinc\texorpdfstring{$^2$}{\^{}2} zero law for squarefree moduli]{The \texorpdfstring{$\sinc^2$}{sinc\^{}2} Zero Law for Squarefree Moduli}

\author{Brayden R.\ Sanders \and M.\ Gish}
\address{7Site LLC, Hot Springs, Arkansas, USA}
\email{brayden@7site.co}

\author{Brayden R.\ Sanders \and M.\ Gish}
\address{Independent Researcher, Hot Springs, Arkansas, USA}
\email{monica.gish1992@gmail.com}

\date{\today}

\keywords{sinc function, squarefree modulus, smallest prime factor,
coprimality partition, divisibility, Fourier coefficients}

\subjclass[2020]{11A41, 11N05, 42A16}

\begin{document}

\begin{abstract}
For any integer $b > 1$ and any positive integer $k$, the squared
sinc function satisfies the divisibility biconditional
$\sinc^2(k/b) = 0 \iff b \mid k$. We re-scope this elementary
identity to the family of squarefree moduli and combine it with the
First-G Event Localization Theorem of \cite{J04} to obtain the
following sharper statement: for a squarefree integer
$b = p_1 p_2 \cdots p_r$ with $p_1 < \cdots < p_r$, the smallest
positive integer $k$ at which $\sinc^2(k/d) = 0$ for at least one
non-trivial divisor $d \mid b$ is exactly $k = p_1$, the smallest
prime factor of $b$. The corridor
$\{1/b, 2/b, \ldots, (b-1)/b\}$ then closes at the prime factors of
$b$ in a layered hierarchy: each prime $p_i \mid b$ first turns on
the $\sinc^2(\cdot/p_i)$-zero at $k = p_i$, and the full
$\sinc^2(\cdot/b)$-zero arrives only at $k = b$. The basic
biconditional is verified by exact arithmetic for all primes
$p \in \{3, 5, 7, \ldots, 199\}$ (the squarefree case with one prime
factor); the verification script
\texttt{proof\_d25\_loop\_closure.py} is supplied as electronic
supplementary material.
\end{abstract}

\maketitle

%% -------- section 1: setup --------
\section{Setup and Statement}\label{sec:setup}

The normalized sinc function is defined as
\begin{equation}\label{eq:sinc-def}
  \sinc(x) =
  \begin{cases}
     \dfrac{\sin(\pi x)}{\pi x} & x \neq 0,\\[4pt]
     1 & x = 0,
  \end{cases}
\end{equation}
and its square $\sinc^2(x) = \sinc(x)^2$. The function appears in
signal processing as the Shannon-sampling kernel
\cite{Shannon1949}, in random matrix theory as the complement of
Montgomery's pair-correlation function for Riemann zeros
\cite{Montgomery1973}, and---as we make precise here---directly in
the elementary divisibility structure of $\Z$.

For an integer $b > 1$ we study the values $\sinc^2(k/b)$ along the
rational corridor $\{1/b, 2/b, \ldots\}$. The basic algebraic fact
is the following divisibility biconditional, which we record first:

\begin{lemma}[Divisibility biconditional]\label{lem:basic}
For any integer $b > 1$ and any positive integer $k$,
\begin{equation}\label{eq:basic}
  \sinc^2(k/b) = 0 \iff b \mid k.
\end{equation}
\end{lemma}

\begin{proof}
Since $k/b \neq 0$ for $k \ge 1$, we have
$\sinc^2(k/b) = \sin^2(\pi k/b) / (\pi k/b)^2$, which vanishes iff
$\sin(\pi k/b) = 0$, iff $\pi k/b \in \pi \Z$, iff $k/b \in \Z$,
iff $b \mid k$.
\end{proof}

Lemma~\ref{lem:basic} is uniform in $b$: it does not distinguish
prime moduli from composite moduli. To extract a statement that
depends genuinely on the prime structure of $b$, we restrict to
\emph{squarefree} $b$ and pass simultaneously over all non-trivial
divisors. This is the content of Section~\ref{sec:main}.

%% -------- section 2: main theorem --------
\section{The Sinc$^2$ Zero Law for Squarefree Moduli}\label{sec:main}

Throughout this section, $b > 1$ is a squarefree integer with
canonical factorization $b = p_1 p_2 \cdots p_r$, where
$p_1 < p_2 < \cdots < p_r$ are the distinct prime factors of $b$.
We write $\spf(b) = p_1$ for the smallest prime factor.

\begin{definition}[Divisor sinc-zero set]
For a squarefree integer $b > 1$ and a positive integer $k$, the
\emph{divisor sinc-zero set} of the corridor position $k$ relative
to $b$ is
\begin{equation}\label{eq:Z}
  Z_k(b) \;=\; \{ d : d \mid b, \ d > 1, \ \sinc^2(k/d) = 0 \}.
\end{equation}
\end{definition}

By Lemma~\ref{lem:basic}, $d \in Z_k(b)$ iff $d \mid b$, $d > 1$,
and $d \mid k$. Equivalently, $Z_k(b) = \{ d : d \mid \gcd(k, b),
\ d > 1 \}$.

\begin{theorem}[Squarefree Sinc$^2$ Zero Law]\label{thm:main}
Let $b > 1$ be squarefree with $\spf(b) = p_1$. Then:
\begin{enumerate}
\item[(i)] $Z_k(b) = \emptyset$ for every $k$ with
$1 \le k < p_1$;
\item[(ii)] $Z_{p_1}(b) = \{p_1\}$, and $p_1$ is the unique
non-trivial divisor of $b$ for which $\sinc^2(p_1 / p_1) = 0$ at
corridor position $k = p_1$.
\end{enumerate}
In particular, the smallest positive integer $k$ at which
$\sinc^2(k/d) = 0$ for at least one non-trivial divisor $d \mid b$
is $k = p_1 = \spf(b)$.
\end{theorem}

\begin{proof}
By Lemma~\ref{lem:basic}, $\sinc^2(k/d) = 0 \iff d \mid k$ for every
$d > 1$. Hence $Z_k(b)$ is exactly the set of non-trivial divisors
of $\gcd(k, b)$.

\emph{Part (i).} Fix $k$ with $1 \le k < p_1$. Suppose for
contradiction that $\gcd(k, b) > 1$; then $\gcd(k, b)$ has a prime
divisor $q$ with $q \mid b$, so $q \ge p_1$, but $q \mid k$ forces
$q \le k < p_1$, a contradiction. Therefore $\gcd(k, b) = 1$, so
$Z_k(b) = \emptyset$.

\emph{Part (ii).} At $k = p_1$, the integer $p_1$ divides itself and
divides $b$, so $p_1 \in Z_{p_1}(b)$. Conversely, suppose
$d \in Z_{p_1}(b)$. Then $d \mid b$, $d > 1$, $d \mid p_1$. Since
$p_1$ is prime, $d \in \{1, p_1\}$, and $d > 1$ forces $d = p_1$.
Hence $Z_{p_1}(b) = \{p_1\}$.

The combined statement follows: parts (i) and (ii) say that the
first $k \ge 1$ for which $Z_k(b) \ne \emptyset$ is $k = p_1$.
\end{proof}

\begin{remark}[Why this is not Lemma~\ref{lem:basic}]
\label{rem:nontrivial}
Lemma~\ref{lem:basic} alone does not separate primes from
composites: it states a uniform fact about every modulus.
Theorem~\ref{thm:main} is strictly stronger because it identifies
the smallest $k$ at which any non-trivial divisor of $b$ produces
a $\sinc^2$ zero, and that $k$ is $\spf(b)$ \emph{precisely
because} $\spf(b)$ is the smallest member of a non-empty set of
primes. Replacing $\spf(b)$ by any composite would break the
contradiction step in Part (i): if $b = 15$ and one tried
$k = 6 < 15$, the divisor $d = 3$ already satisfies $d \mid k$, so
the corresponding $Z_6(15)$ already contains $3$; the localization
fails. The squarefree hypothesis is needed only to ensure that the
ladder of prime divisors of $b$ has a clean smallest-element
structure (every prime power
$p^a \mid b$ collapses to its prime base).
\end{remark}

\begin{remark}[Relation to the First-G Law]\label{rem:first-g}
Theorem~\ref{thm:main} is the $\sinc^2$ shadow of the First-G Event
Localization Theorem of \cite{J04}, which states that for every
$b > 1$ with smallest prime factor $p_1$, the first non-coprime
element of the alphabet $\{1, \ldots, k\}$ relative to $b$ appears
at exactly $k = p_1$. The two statements are equivalent under the
identification $\sinc^2(k/d) = 0 \iff d \mid k \iff d \in
\{\text{prime divisors of } \gcd(k, b)\}$. We refer the reader to
\cite{J04} for the algebraic version and full corollary structure
(stability windows, prime-indexed phase transitions, CRT idempotent
count, instability ranking of primes).
\end{remark}

%% -------- section 3: corollaries --------
\section{Corollaries}\label{sec:cor}

The following three corollaries record familiar consequences of
Theorem~\ref{thm:main} that arise naturally in the corridor
picture. They are equivalent to the analogous corollaries in the
First-G companion \cite{J04}; we restate them in $\sinc^2$ language
for the convenience of the spectral reader.

\begin{corollary}[Layered loop closure]\label{cor:loop}
Let $b = p_1 p_2 \cdots p_r$ be squarefree. As $k$ increases from
$1$ to $b$, the divisor sinc-zero set $Z_k(b)$ grows by inclusion in
exactly $\tau(b) - 1$ stages, where $\tau(b) = 2^r$ is the divisor
count: stage $j$ adds the non-trivial divisors of $b$ whose smallest
prime factor first divides $k$ at $k = p_{j}$. The corridor
$\{1/b, \ldots, b/b\}$ closes (in the sense $Z_b(b) = \{
\text{all non-trivial divisors of } b\}$) at $k = b$.
\end{corollary}

\begin{proof}
By Lemma~\ref{lem:basic}, $d \in Z_k(b)$ iff $d \mid \gcd(k, b)$ and
$d > 1$. The map $k \mapsto \gcd(k, b)$ is monotone non-decreasing
along the divisor lattice of $b$ as $k$ increases by $1$ at a time,
and changes value exactly when $k$ acquires a new prime factor of
$b$. Each such acquisition first happens at some $k = p_j$ for some
$j \le r$. At $k = b$, $\gcd(b, b) = b$, so $Z_b(b)$ contains every
non-trivial divisor of $b$.
\end{proof}

\begin{corollary}[Prime-indexed amplitude transitions]
\label{cor:prime-indexed}
The amplitudes $\sinc^2(k/b)$ along the corridor
$k \in \{1, 2, \ldots, b-1\}$ are strictly positive
\emph{everywhere} (since $b \nmid k$ in this range), but the
corridor amplitudes $\sinc^2(k/p_j)$ for the divisors
$p_j \mid b$ each cross zero for the first time at $k = p_j$.
The set of corridor positions at which any such divisor amplitude
first vanishes is exactly $\{p_1, p_2, \ldots, p_r\}$, the prime
factors of $b$.
\end{corollary}

\begin{proof}
Direct application of Lemma~\ref{lem:basic} to each divisor of $b$
in turn: $\sinc^2(k/p_j) = 0 \iff p_j \mid k$, and the smallest
$k \ge 1$ for which $p_j \mid k$ is $k = p_j$.
\end{proof}

\begin{corollary}[Stability window]\label{cor:stab}
The interval $\{1/b, 2/b, \ldots, (p_1 - 1)/b\}$ is sinc-zero free
in the strong sense that $\sinc^2(k/d) > 0$ for every
$k \in \{1, \ldots, p_1 - 1\}$ and every non-trivial divisor
$d \mid b$. The width of this stability window is $p_1 - 1$, and
it depends only on $\spf(b)$, not on the larger prime factors
$p_2, \ldots, p_r$.
\end{corollary}

\begin{proof}
Part (i) of Theorem~\ref{thm:main} says $Z_k(b) = \emptyset$ for
$k < p_1$, which is exactly the strong stability statement.
Width and uniformity in $\spf(b)$ are immediate.
\end{proof}

%% -------- section 4: boundary value --------
\section{The Boundary Value $\sinc^2(1/2) = 4/\pi^2$}\label{sec:boundary}

The amplitude at corridor midpoint $k/b = 1/2$ is independent of
$b$ and admits a closed form:
\begin{equation}\label{eq:fourpi2}
  \sinc^2(1/2) = \left( \frac{\sin(\pi/2)}{\pi/2} \right)^{\!2}
              = \left( \frac{2}{\pi} \right)^{\!2}
              = \frac{4}{\pi^2} \approx 0.4053.
\end{equation}
This is a transcendental number, the amplitude of the first
sidelobe of a rectangular spectral gate (a classical fact from
signal processing). It appears here as a derived constant of the
corridor: any corridor of even width $2m$ has its midpoint at
$k = m$, and the corresponding $\sinc^2$ amplitude is exactly
$4/\pi^2$.

Montgomery's pair-correlation function for Riemann zeros is
$R_2(u) = 1 - \sinc^2(u)$ \cite{Montgomery1973}. The complement
$\sinc^2(u)$ and $R_2(u)$ sum to unity, giving a complete spectral
partition. The boundary $4/\pi^2 = \sinc^2(1/2)$ appears in both
frameworks, derived rather than imposed.

%% -------- section 5: companion --------
\section{Connection to the First-G Companion}\label{sec:firstG}

Theorem~\ref{thm:main} is the $\sinc^2$ image of the algebraic
statement proved in the First-G companion paper \cite{J04}:

\begin{quote}
\emph{First-G Event Localization Theorem \cite{J04}.} For every
$b > 1$ with smallest prime factor $p_1$, the smallest $k$ for
which the alphabet $\{1, \ldots, k\}$ contains an element sharing
a prime factor with $b$ is $k = p_1$, and this element is $p_1$
itself.
\end{quote}

The translation is
$\sinc^2(k/d) = 0 \iff d \mid k \iff d \mid \gcd(k, b)$, so the
non-trivial divisor in the First-G formulation is exactly the
divisor recorded in our $Z_k(b)$. The two papers share a
verification script: the squarefree case verifies both
formulations simultaneously.

The continuum-limit reading of \cite{J04} gives the harmonic
pre-echo function
\begin{equation}\label{eq:R}
  R(k, b) \;=\; \frac{\sin^2(\pi k / b)}{k^2 \, \sin^2(\pi/b)}
  \;\longrightarrow\; \sinc^2(k/b) \quad (b \to \infty),
\end{equation}
which recovers the $\sinc^2$ field exactly. The two results are
two faces of the same algebraic structure---one discrete (First-G),
one continuous ($\sinc^2$ corridor).

%% -------- section 6: verification --------
\section{Verification}\label{sec:verify}

Theorem~\ref{thm:main} is verified for the squarefree case
$b = p$ prime by the script
\texttt{proof\_d25\_loop\_closure.py}, which checks for all primes
$p \in \{3, 5, 7, 11, \ldots, 199\}$ (46 primes) that:
\begin{itemize}\setlength\itemsep{1pt}
  \item $\sinc^2(k/p) > 0$ for all $k \in \{1, \ldots, p-1\}$
        (exact rational arithmetic via \texttt{fractions.Fraction}
        and \texttt{sympy} on the rational input; floating-point
        $\sinc^2$ on the real output);
  \item $\sinc^2(p/p) = \sinc^2(1) = 0$ exactly;
  \item zero exceptions across all 46 primes.
\end{itemize}
The general squarefree case $b = p_1 p_2 \cdots p_r$ is verified by
the companion script \texttt{proof\_first\_g\_event.py}
\cite{J04proof}, which checks all squarefree $b \le 500$ (153
moduli, 36{,}662 $(b, k)$ pairs) with zero exceptions. We omit a
separate enumeration here because the squarefree-multiprime case
follows from the First-G companion via the equivalence stated in
Section~\ref{sec:firstG}.

The script \texttt{proof\_d25\_loop\_closure.py} (approximately
260 lines) is supplied as electronic supplementary material and is
also archived at DOI~\texttt{10.5281/zenodo.18852047}.

%% -------- section 7: limits --------
\section{What This Result Does Not Claim}\label{sec:limits}

This paper does not claim:
\begin{itemize}\setlength\itemsep{1pt}
  \item a new proof of the infinitude of primes;
  \item a connection to the Riemann Hypothesis beyond the
        Montgomery-pair-correlation observation in
        Section~\ref{sec:boundary};
  \item that the fold position
        $\sinc^2(x^*) = 1/2$, with $x^* \in (0,1)$, is computable
        in closed form;
  \item any result about the distribution of primes;
  \item that the squarefree hypothesis is essential. (The
        non-squarefree case follows from the squarefree case
        applied to the radical $\mathrm{rad}(b)$ of $b$, but
        introduces multiplicities that are tangential to the
        present statement.)
\end{itemize}
Theorem~\ref{thm:main} is a finite, elementary algebraic fact
about $\sinc^2$ evaluated at rational points whose denominator
ranges over the divisors of a squarefree integer.

%% -------- references --------
\begin{thebibliography}{99}

\bibitem{J04}
B.~R.~Sanders and M.~Gish.
\newblock The First-G Event in the Coprimality Partition:
Stability Windows, CRT Idempotent Count, and Prime-Indexed Phase
Transitions.
\newblock Submitted to \emph{Integers}, 2026 (companion paper to
the present manuscript).

\bibitem{J04proof}
B.~R.~Sanders and M.~Gish.
\newblock \texttt{proof\_first\_g\_event.py}: exhaustive
verification of the First-G Event Localization Theorem for all
squarefree $b \le 500$.
\newblock Electronic supplementary material to \cite{J04}, 2026.

\bibitem{Montgomery1973}
H.~L.~Montgomery.
\newblock The pair correlation of zeros of the zeta function.
\newblock In \emph{Analytic Number Theory}, Proceedings of Symposia
in Pure Mathematics, volume~24, pages 181--193. American
Mathematical Society, Providence, RI, 1973.

\bibitem{Shannon1949}
C.~E.~Shannon.
\newblock Communication in the presence of noise.
\newblock \emph{Proceedings of the IRE}, 37(1):10--21, 1949.

\end{thebibliography}

%% -------- end matter --------
\vfill
\noindent
\footnotesize{%
\textsc{Provenance.}\; Source manuscript
\texttt{WP\_SINC2\_ZERO\_LAW.md} in the 7Site Research repository
(DOI~\texttt{10.5281/zenodo.18852047}). Re-scoped 2026-05-07 from
the 2026-04-18 prime-only draft to the squarefree-modulus
formulation, with the prime-specific content routed through the
First-G companion paper \cite{J04}. Verification script
\texttt{proof\_d25\_loop\_closure.py} runs in standard CPython with
no external dependencies; runtime is under five seconds for the
46-prime sweep.}

\end{document}
